% ===================================================================
% The answer-equivalent: definition, protocol, price, mechanism, boundary.
% Included from main.tex after the audit sections, which supply the
% design rationale for every choice made here.
% ===================================================================

\section{The answer-equivalent}
\label{sec:equiv}

\paragraph{Definition.} Let $S$ be a strictly proper ordinal score --- we use the
ranked probability score, normalised by the $K{-}1$ thresholds, so $S$ is
comparable across the $5$- and $7$-point inventories we study. Let a respondent
$i$ have observed answers $x_i$ and held-out answers $y_i$ on disjoint item sets.
Write $D_m$ for a conditional model fitted on training respondents that sees $m$
of respondent $i$'s observed answers, and $\widehat{F}$ for the simulator under
test. The answer-equivalent is
\[
m^{*}(\widehat{F})=\min\{m: \mathbb{E}\,S(D_m, y)\le \mathbb{E}\,S(\widehat{F}, y)\}.
\]
We report both this integer crossing, which is what a practitioner acts on, and
its linear interpolation between the bracketing budgets, which is the honest
effect size: the crossing is coarse and rounds a simulator up by as much as a
whole answer.

\paragraph{What the simulator is allowed.} $\widehat{F}$ is the persona readout
after population calibration; the raw readout scores worse than the population
prior, so pricing it raw would price the prompt format rather than the model. We
fit $\widehat{F}_{ij}=\Pi_{\mathrm{CDF}}[F_{0j}+\alpha(Q_{ij}-\bar{Q}_j)]$, with
$Q$ the model's own CDF from its digit logprobs, $F_0$ a population CDF fitted on
training respondents, $\bar{Q}$ its mean CDF on a calibration pool,
$\Pi_{\mathrm{CDF}}$ the exact $L_2$ projection onto monotone CDFs, and $\alpha$
chosen on calibration respondents. This is deliberately the most favourable
treatment: it hands the simulator a correctly located population distribution for
free and asks only whether its \emph{person-specific} deviations are worth
anything.

\paragraph{What the baseline is allowed.} $D_m$ is a ridge over the one-hot
observed answers onto the threshold indicators, projected onto monotone CDFs.
Its penalty is retuned at \emph{every} budget $m$, on the same calibration
respondents that select $\alpha$. This is not a detail: holding the penalty fixed
overfits the small budgets, and in development that alone manufactured an
apparent $15\%$ relative gain for the simulator at $m{=}100$ which vanished once
the baseline was tuned at each budget. A comparison across label budgets that
does not retune baseline capacity is measuring the baseline's misspecification.

\paragraph{Protocol.} Conditioning and target items are disjoint by
construction; respondents are split into a training block, a calibration half on
which every hyperparameter is selected, and an evaluation half on which every
reported number is computed frozen; uncertainty is a paired bootstrap over
evaluation respondents. Appendix~\ref{app:protocol} gives the full specification,
including the reverse-orientation correction, the common-support rule and the
multiplicity control.

\section{What a persona is worth}
\label{sec:price}

We price $89$ cells: every model in our artifact set, on IPIP-NEO-120, SD3,
IPIP-FFM-50 and HEXACO, on two disjoint respondent panels, plus the IPIP-NEO-300
unseen-item holdout of \S\ref{sec:transfer}.

\paragraph{The instrument is stable.} Across the $32$ (model, instrument) pairs
scored on both panels, $m^{*}$ on the \texttt{val} panel and on the \texttt{test2}
panel correlate at Spearman $\rho{=}0.78$ ($p{=}1.7\times10^{-7}$). The price is a
property of the model and instrument, not of a panel draw. Without this, nothing
below would be worth reporting.

\paragraph{The price.} Table~\ref{tab:ladder} gives the IPIP-NEO-120 ladder.
Median $m^{*}$ per inventory is $0.00$ on HEXACO, $0.06$ on IPIP-FFM-50, $0.32$
on IPIP-NEO-120, $0.40$ on SD3 and $0.61$ on the IPIP-NEO-300 holdout
(App.~\ref{app:equivfull}). The strongest persona we
measure on the short form is worth $3.07$ of $60$ answers; the weakest is worth
exactly zero, in the strong sense that calibration selects $\alpha{=}0$ and
discards its person-specific deviations entirely. Across all $89$ cells the median
is $0.19$ answers.

\begin{table}[t]\centering\footnotesize
\begin{tabular}{lrrr}
\toprule
Model & raw $S$ & $\alpha^{*}$ & $m^{*}$ (of 60)\\
\midrule
Qwen3.6-35B-A3B & 0.172 & 0.85 & \textbf{3.07}\\
Qwen2.5-32B & 0.213 & 0.30 & \textbf{2.15}\\
Qwen2.5-72B & 0.202 & 0.45 & \textbf{1.66}\\
Qwen3.8-27B & 0.183 & 1.10 & \textbf{1.31}\\
QwQ-32B-Preview & 0.202 & 0.70 & \textbf{1.29}\\
granite-4.2-30b & 0.243 & 0.45 & \textbf{0.45}\\
gemma-4-12B-it & 0.267 & 0.35 & \textbf{0.37}\\
Mistral-Small-24B-2501 & 0.324 & 1.00 & \textbf{0.33}\\
\midrule \multicolumn{4}{c}{\itshape 5 further models, App.~\ref{app:equivfull}}\\ \midrule
granite-4.1-8b & 0.356 & 0.20 & \textbf{0.02}\\
Apertus-8B-2509 & 0.286 & 0.00 & \textbf{0.00}\\
\bottomrule
\end{tabular}
\caption{Answer-equivalent on IPIP-NEO-120, averaged over both panels. $S$ is the
normalised ordinal RPS of the raw readout. $\alpha^{*}$ is the selected
calibration coefficient. Raw accuracy does not order the ladder;
$\alpha^{*}$ does.}
\label{tab:ladder}
\end{table}

\paragraph{Reading the number.} $3.07$ of $60$ is a real, replicable, non-zero
price, and it is small. Both halves matter. A practitioner with a $60$-item
instrument and a budget for three more questions should ask them. A practitioner
choosing between open models should read the ladder, because the spread across
models --- $0$ to $3.07$ --- is larger than any effect reported for prompt
optimisation on these systems.

\subsection{Transfer to an unseen item family}
\label{sec:transfer}

The splits above put every target facet's siblings in the conditioning half. The
IPIP-NEO-300 holdout removes that crutch: $300$ respondents answered the full
inventory, $120$ of whose items are exactly the short form used everywhere else,
and the other $180$ --- spanning all $30$ facets --- are items the same person
answered that no arm has been fitted on, with only $150$ training respondents
behind them. This is the thinnest supervision we measure and the most favourable
regime for a simulator, and it is where the price is highest: $\mathbf{5.38}$ of
$120$ answers, against $3.24$ of $60$ on the short form, with the model ordering
preserved. A persona is worth most exactly where human labels are scarcest.

Presentation matters as much as quantity. Given its persona as $m$ literal (item,
answer) pairs rather than a $30$-line profile, the same model prices at $0.00$,
$0.19$, $0.49$, $1.10$ for $m{=}5,20,60,120$ --- monotone, as it must be --- while
the profile built from those same $120$ answers is worth almost five times the
transcript of all of them ($\alpha^{*}$ $0.9$ against $0.3$). The direction is not
universal: Apertus-0.5B agrees ($0.13$ against $0.06$) but gemma-4-E2B reverses it
($0.22$ against $0.54$). Three models cannot settle which presentation wins; what
the price shows is that the choice is worth several observations and belongs in
the report (App.~\ref{app:transfer}).

\section{Why some personas are worth more}
\label{sec:why}

The ladder is not ordered by accuracy. The raw score of the best-priced model
($0.172$) and of \texttt{gpt-oss-20b} ($0.188$) are close, but their prices differ
by more than a factor of ten. Across all $89$ cells, raw score correlates with
$m^{*}$ at $\rho{=}-0.48$, and the \emph{calibrated} score --- the thing one would
naively report --- does not correlate with it at all.

What does order the ladder is $\alpha^{*}$, the share of a model's person-specific
deviation that survives calibration: $\rho{=}+0.72$ ($p{<}10^{-14}$), and
$\rho{=}+0.66$ after partialling out raw score. A model whose forecast barely
moves between people offers calibration nothing to keep; \texttt{Apertus-8B} has
$\alpha^{*}{=}0$ and a price of exactly zero, despite a middling raw score.

$\alpha^{*}$ is fitted against the same score $m^{*}$ is derived from, so on its
own it is a tight proxy rather than an independent cause. We therefore measured
the belief tensor directly, with no fitting and no labels. Controlling for raw
score, the \emph{entropy} of the model's person-conditional belief predicts the
price at $\rho{=}-0.52$ ($p{=}2.4\times10^{-7}$, $n{=}88$) while its marginal
correlation is null ($\rho{=}-0.03$): sharper conditional beliefs are worth more
real answers, independently of accuracy. Each step of the mechanism is then
measured rather than assumed --- sharper beliefs, larger person-specific
deviations, more of them surviving calibration, a higher price.

Between-person \emph{spread} of the belief does not predict the price
($\rho{=}+0.13$, $p{=}0.22$), nor do the cluster-level variance ratios of
\S\ref{sec:mechanism}. A population-level dispersion audit does not tell you what
a persona is worth to an individual; it is the sharpness of the conditional
belief, not the spread of the mixture, that carries the value.

\section{The boundary: personas are redundant given real answers}
\label{sec:boundary}

The price is quoted against a model given $m$ answers. The complementary question
is whether a persona adds anything to a model that already has \emph{all} of the
respondent's observed answers. It does not.

We screen $160$ cells. For each, three combination families are fitted on the
calibration half and frozen on the evaluation half: a convex weight between the
conditional model and the calibrated persona; a linear combiner taking the
persona's CDF as additional features beside the conditional model's; and a
gradient-boosted combiner over the same features. The control in each case is the
identical combiner, of identical capacity, fitted on identical rows, without the
persona features --- so the comparison isolates the persona and nothing else.

% eight columns do not fit one ACL column: as \begin{table} this overflowed the
% column by 98pt (about 1.4in of text in the margin). table* spans both.
\begin{table*}[t]\centering\footnotesize
\begin{tabular}{lrrrrrrr}
\toprule
Corpus & cells & LLM & prior & decoder & $\bar w^{*}$ & sig.\ cells & best $\Delta$\\
\midrule
HEXACO & 40 & 0.222 & 0.159 & 0.119 & 0.040 & 4/80 & +0.0001\\
IPIP-FFM-50 & 24 & 0.239 & 0.166 & 0.121 & 0.000 & 0/48 & +0.0000\\
IPIP-NEO-120 & 54 & 0.247 & 0.156 & 0.110 & 0.001 & 2/108 & +0.0000\\
SD3 & 42 & 0.215 & 0.167 & 0.130 & 0.000 & 2/84 & +0.0001\\
\bottomrule
\end{tabular}
\caption{Incremental value of a persona given the respondent's own observed
answers. $\bar w^{*}$ is the mean selected convex weight on the persona.
Significant cells are \emph{uncorrected} counts, shown to make the multiplicity
visible: across all $426$ one-sided tests, one survives Benjamini--Hochberg at
$q<0.05$, and it reverses sign on its own disjoint panel.}
\label{tab:incremental}
\end{table*}

The selected convex weight on the persona is $0$ in $127$ of $160$ cells, $0.05$
in the remaining $33$, and never higher. The linear combiner finds a nominally
significant positive increment in $0$ of $160$ cells. Across the full grid ---
$426$ one-sided tests --- $35$ are nominally significant at $5\%$, which is what
chance produces, and exactly \emph{one} survives Benjamini--Hochberg at $q<0.05$:
Qwen3.8-27B on HEXACO, \texttt{val} panel, gradient-boosted combiner,
$+0.00092$ $[+0.00046,+0.00140]$, $q{=}0.026$.

That cell does not survive scrutiny: on the disjoint \texttt{test2} panel the gain
reverses sign ($-0.00034$), and at four times the panel size it falls to
$+0.00006$ --- a small-sample artefact, as is every other nominally significant
HEXACO cell, all of which sit on the same $64$-respondent panel. We report it
rather than filter it out; a boundary claim that discards its one exception is
not a boundary claim.

This is a null reported as a null, not a failed superiority test: the combiner
gets every advantage its control gets, and it is the selected weight --- zero in
$127$ of $160$ cells --- rather than a $p$-value that concedes.

This is the sharp statement the price makes precise. A persona carries a small,
real, model-scaling amount of personal information --- and a handful of the
person's actual answers renders it redundant.
